\section{Introduction}\label{sec:intro}

\subsection{The problem}

Can a deterministic multitape Turing machine transpose a $d\times d$ matrix with $\clog d$-bit entries in time linear
in the input length $L=d^2\clog d$? We call the affirmative statement~(T); \cref{sec:question} states it precisely. The
folklore algorithm takes $O(L\log d)$ steps. Transposition is the data-movement step of fast integer multiplication on
tapes, and by a reduction of Harvey and van der Hoeven~\cite{HvdH2025} a linear-time multiplier would give~(T). So a
proof that (T) fails would be the first superlinear lower bound for off-line multiplication on multitape machines, and a
proof that it holds would be a new way of moving data.

Sixteen papers and eight memos of this series have studied the question since October~2026. They contain reductions,
restricted impossibility theorems, constructions of codes and many formal proofs. They do not contain the answer, and a
reader who starts today has to read all of them to learn where the difficulty is. This paper is written for that reader.

\subsection{What this paper does}

It does three things.

\paragraph{A synthesis (\cref{sec:series}).} We state the main results of the series in one place, by theme: the
reductions, what head motion must do, depth reduction, what is proved about storage codes, what is false about them, and
adaptive machines. Nothing is reproved. Every statement is cited, with its main hypotheses.

\paragraph{A diagnosis (\cref{sec:wall}).} For oblivious machines the series has reduced~(T) to one function: the best
rate $r(\eps)$ of a storage code on the bounded-depth cone relations that remain of linear-time schedules when a genie
reveals $\eps L$ bits. A transposer needs $r(\eps)\ge1-\eps$ for every~$\eps$. The matching of each stored symbol with
one partner gives $r\ge\tfrac12$, and whether $r$ exceeds $\tfrac12$ at all on the relations of real machines is open in
both directions. Three kinds of argument have reached known limits:
\begin{itemize}[leftmargin=1.5em]
\item conditions on motion that ignore the data, because a machine that transposes nothing satisfies the cut
inequalities behind them; conditions that do see the outputs are single-scale, and one cut serves every pair of row
sets;
\item bounds for codes that ignore the machine, because abstract relations with the two properties used so far,
bounded degree and paired-distance separation, carry codes of rate tending to one;
\item counting passes, because it saturates one scale at a time.
\end{itemize}
They stop for one reason. The logarithm of the folklore bound is a sum over $\log_2d$ address scales, each costing $L$,
and every tool of the series that applies to all machines looks at one scale or at none; the logarithm is proved only
where the coding or the geometry is tied down. The missing statement has to tie what is known to where it can be, for
all schedules, and the coding properties proved so far are properties of the cone \emph{relation}, which forgets that
neighbouring outputs of a machine are produced by the same rule.

\paragraph{Two theorems (\cref{sec:order,sec:symmetry}).} We prove two rigidity theorems for the swap in a model where
the code and its decoder are visible together. Let $V=F^d$ over any field, let $\sw$ be the swap on $V\otimes V$, that
is, $X\mapsto X^{\tr}$ on matrices, and let $K=\sum_aA_a\otimes B_a$ act by $X\mapsto\sum_aA_aXB_a^{\tr}$. The code
$\ker(\sw-K)$ is the set of matrices that $K$ transposes correctly. The symmetric matrices, for $K=1$, have dimension
$d(d+1)/2$.

\begin{itemize}[leftmargin=1.5em]
\item \emph{Order (Theorem~FL, \cref{thm:FL}).} For every complete flag $W_1\subset\dots\subset W_d$ of $V$, with
$g(i)=\dim\sum_aB_a(W_i)$,
\[
  \dim\ker(\sw-K)\le\sum_{i=1}^dg(i),
\]
whatever the $A_a$ are. So two Kronecker terms never beat the symmetric matrices, over any field; neither do commuting
$B_a$, for instance any convolutions of any number of rows; and neither does any combination of rotations, reflections,
dilations and decimations when $d=p^m$ in characteristic~$p$. The bound is attained for every profile~$g$.
\item \emph{Symmetry (Theorem~TP, \cref{thm:TP}).} For any three operators on an $N$-dimensional space,
$\dim\ker(\sw-K)\le\frac12(N+\dim\ker[\sw,T]+\rank[K,T])$, an elementary lemma. For the swap, $\dim\ker[\sw,T]\ge d$
always. If row $c$ of $X$ is a fixed linear function of $X$ with its columns rotated by $c$, the same function for all
$c$ in each of $M$ intervals, then
\[
  \dim\{X\}\le\frac{d^2+Md}2 ,
\]
on a cycle and on a line, with equality for equal intervals: there the largest code has the dimension of the coarse
$M\times M$ block transposition. For $M$ equal intervals flags give only rate $1-1/(2M)$. Variants bound decoders with $s$
modulation frequencies by rate $s/(s+1)$ and show that every address digit must be treated asymmetrically. The
hypothesis on the classes is needed: for $d=6$, two classes with no almost-period carry a code of dimension $25>24$
(\cref{prop:aperiodic}), found by a SAT solver and checked in Lean, and for every fixed number $M\ge2$ of classes the gain over
the bound is unbounded (\cref{sec:two-class}, added after a memo to this paper).
\end{itemize}

In words: to beat the pairing value, the decoder may have neither a hidden order nor a symmetry, on either index, that
the swap does not share. The second theorem uses exactly the feature of machines that cone relations forget.

\paragraph{What the theorems are not.} They are statements about the linear one-hop model. They are not lower bounds
for machines, and (T) remains open. \Cref{sec:construction} lists what they exclude and the three steps that would be
needed to carry them to machines; \cref{sec:small} gives exact values at small sizes; \cref{sec:open} states the
problems.

\paragraph{A subsequent memo, and the revisions of this paper.} After revision \texttt{b2c389a} of this paper in its
repository, a memo~\cite{ClassUniformMemo2026} proved, in Lean, that the two classes of \cref{prop:aperiodic} carry a
code of dimension $25$ over every field in which $3\ne0$; that class-uniform codes can be blown up, which gives gains
that grow with the number of classes; and bounds for the cost of advice in the one-hop model. Revision
\texttt{14b157a} of this paper cited the memo where it bears on the text, and corrected three passages:
\begin{itemize}[leftmargin=1.5em]
\item the charge for a bit that a genie reveals (\cref{sec:to-machines,prob:machines}): the earlier text charged every
revealed bit as a change of regime, at half a row, where advice that enters additively changes no decoder and half a
row is the price of a change that affects a whole row;
\item \cref{prob:aperiodic}(c), which the memo answered as it was worded and which was then stated for a fixed number
of classes;
\item the first sentence of \cref{sec:would-need}, which had dropped a hypothesis of the result it summarizes.
\end{itemize}
It added \cref{prop:entry,cor:entry-advice}, which state the corrected charge, with formal proofs of their bounds in
terms of the number of changes, and it updated \cref{rem:small}(d) and \cref{prob:aperiodic}(b).

An update of the memo then answered that form of \cref{prob:aperiodic}(c) for two classes: for every $d=4k\ge8$ it
gives an explicit rule, over fields in which $2\ne0$, whose code exceeds $(d^2+2d)/2$ by $d/4-1$, so that the gain of
two classes is unbounded. Its proof is written in the memo, which also proves in Lean, for $d=8$, $12$ and~$16$, that
the two classes carry a code of that dimension. It leaves open characteristic~$2$ and three or more classes.
Revision \texttt{ac70dc4} of this paper cited that update and added \cref{sec:two-class}, which settles both.
\begin{itemize}[leftmargin=1.5em]
\item \Cref{prop:two-class}: the memo's rule works over every field, $\F_2$ included, once its coefficients $2$ and
$\tfrac12$ are replaced by a nonzero $\lambda$ and its inverse.
\item \Cref{lem:mixed-blowup,cor:every-M}: a blow-up with a code of its own at each position turns the two-class codes
into codes for any fixed number $M\ge2$ of classes, with a gain of order $Md/8$ over $(d^2+Md)/2$. So the gain is
unbounded for every fixed number $M\ge2$ of classes, over every field.
\end{itemize}
These three results are proved in the text and are not formalized. \Cref{prob:aperiodic}(c) now asks for the rate of
growth.

Two papers of the series followed revision \texttt{ac70dc4}. One is about the original runs of adaptive
machines~\cite{AdaptiveTransfer2026}. The other is about the class-uniform rules of \cref{sec:symmetry}, their
codes and their composition~\cite{SharedDecoders2026}. Revision \texttt{3911d87} cites both where they bear on the text,
and changes the following.
\begin{itemize}[leftmargin=1.5em]
\item \emph{Bounds for classes of any shape} (\cref{prob:aperiodic}(a), \cref{sec:would-need}). For two classes the
second of these papers proves $\tfrac{41}{64}d^2+O(d^{3/2})$, where this paper had $\tfrac34d^2+O(d)$, and for $M$
classes a deficit of order $M^{-1/2}$, where this paper had $1/(2M)$. For partitions with a period, for which
\cref{thm:TP} shows that the gain is at most linear in~$d$, it determines the largest code of a lift under
hypotheses that are stated with \cref{prob:aperiodic}(c).
\item \emph{Which linearity the bounds need} (\cref{sec:sets}, new). Theorem~TP needs the code to be a subspace, but
not the decoders to be linear. For sets of matrices that are not subspaces it fails completely: one decoder for all
rows is correct on a fraction $1/(2d)$ of all matrices. Both facts are from that paper. That revision adds
\cref{prop:sets}: the elimination step of Theorem~FL and the bound for rules that copy one entry hold for sets and
arbitrary functions, as instances of the acyclic-subgraph bound of Bar-Yossef, Birk, Jayram and
Kol~\cite{BBJK2011}; and a remark on marker rows, by which a set of matrices with one function for all rows exceeds
the value $(d^2+d)/2$ of subspaces. One sentence of \cref{sec:symmetry-says} and one of \cref{sec:reading}, which said that
Theorem~TP needs linear decoders, are corrected, and a sentence is added to \cref{sec:symmetry-says}.
\item \emph{The step to machines} (\cref{sec:to-machines}, \cref{prob:machines}). The earlier text took the form of
its question from a count for affine rules. For decoders that are not affine and see the whole matrix there is no such
count. What machines add is that the cones a genie leaves are of bounded size, and whether shared decoders of
bounded window bound a set of matrices is the new \cref{prob:machines}(b). \Cref{sec:to-machines} also cites the
first bound for two hops, and notes that a perfect shuffle of the columns has $d-1$ classes; \cref{sec:fail} notes
what \cref{prop:sets} gives for synchronous access.
\item \emph{Adaptive machines} (\cref{sec:adaptive}, \cref{sec:one-function}, \cref{sec:exhaustions}). The first of
the two papers proves on the machine's own runs a statement of the form that \cref{sec:adaptive} reported as open,
and gives the function of \cref{sec:one-function} a form on those runs.
\item \Cref{prob:unequal} is now said to be about linear decoders.
\end{itemize}
\Cref{prop:sets} is proved in the text and is not formalized. The other results of this paper and their formal proofs
are unchanged, and so are their numbers.

That revision, \texttt{3911d87}, asked \cref{prob:machines}(b) for every window. The revision after it,
\texttt{a445cd8}, answered part of it, in two propositions added to \cref{sec:sets}.
\begin{itemize}[leftmargin=1.5em]
\item \Cref{prop:order}: the acyclic-subgraph bound, applied through an order of the rows or of the columns of the
matrix. For window two it gives at most $q^{d(d+1)/2+d-1}$ matrices, for all maps, all functions and all sets,
which is attained at $d=3$ and is within $d/2$ of an example at every even~$d$; and
for every window it gives $q^{(1/2+o(1))d^2}$ when the offsets, or the differences of the rows from a permutation,
take boundedly many values. Both hold for functions of whole rows, respectively whole columns.
\item \Cref{prop:mutual}: for every $\eps>0$, all but $\eps d^2$ of the pairs of positions that help each other are
related by one of boundedly many translates of the swap, so a code that uses only such pairs stays at the pairing
rate; and $o(d^2)$ positions meet all transitive triangles, which makes the theorems of the seventh paper for
relations without them available on one-class relations.
\item \Cref{prob:machines}(b) is now asked for windows of at least three, and \cref{sec:sets} reports computations
which indicate that for large windows the acyclic-subgraph bound cannot answer it.
\end{itemize}
The two propositions are proved in the text and are not formalized. No earlier result changes, and no number of a
result or of a table.

A further revision adds four propositions after them, on operators whose rows are arbitrary functions of whole rows
of their argument, and \cref{prob:machines}(c).
\begin{itemize}[leftmargin=1.5em]
\item \Cref{prop:rows}: such an operator transposes at most $q^{\sigma}$ matrices, where $\sigma$ is the least sum,
over the orders of the rows, of the place of the last row that each row of the output reads; this is exactly the
acyclic-subgraph bound of the relation, it equals $(d^2+sd)/2$ when each row of the output reads the block of $s$
rows that contains its image under a bijection, and it is at most $wd(d+1)/(w+1)$ when $w$ rows are read.
\item \Cref{prop:stages}: in a pipeline of stages that read rows inside cosets of the address space, with arbitrary
functions and alphabets between the stages, the dimensions of the cosets add up, and $t$ moved digits give at most
$q^{(d^2+2^td)/2}$ transposed matrices, which partial transposes attain. So a set of rate $\tfrac12+\delta$ needs
$\log_2d-\log_2(1/(2\delta))$ moved digits, for every code. Scans with a buffer of one row enter through advice,
and pipelines that mix differently arranged copies of the array through the number of ones of the address, with
$\sqrt{\log d}$ in place of $\log d$.
\item \Cref{prop:expanding}: if the rows read expand, the bound of \cref{prop:rows} is at least
$(\tfrac12+c)d^2$, and by the monotone expanders of Bourgain and Yehudayoff, in the form with three matchings of
Dujmovi\'c, Sidiropoulos and Wood, there are such relations made of the identity and three increasing maps, which
three forward passes gather when the strings that mark the maps are given. So the acyclic-subgraph bound cannot show
that operators of boundedly many such passes transpose only $q^{(1/2+o(1))d^2}$ matrices.
\item \Cref{prop:pieces}: products of codes on boundedly many positions stay at the pairing rate on every such
relation of bounded degree.
\end{itemize}
These are proved in the text and are not formalized, except that \cref{prop:expanding}(b) applies a cited theorem,
of which we use the statement and have not checked the proof. No earlier result changes, and no number of a result
or of a table.

\subsection{How to read this paper and its formal proofs}

The synthesis and the diagnosis contain no numbered results; they cite. The numbered results of
\cref{sec:order,sec:symmetry} each carry a status line.
\begin{itemize}[leftmargin=1.5em]
\item ``Lean:'' names the formal counterpart in the library \leanfile{lean/SwapRigidity/}; the Lean statement is
authoritative, and the text gives an outline.
\item ``Lean (partial):'' names what is formalized and says what is not.
\item ``Not formalized'' is followed by a complete proof.
\end{itemize}
Formalized in Lean~4 with Mathlib, on the standard axioms only: the coordinate form of Theorem~FL, its attainment and
the corollary for a pair of flags; the pigeonhole lemma, its composition property and the lower bound
$\dim\ker[\sw,T]\ge d$; Theorem~TP on a cycle with an arbitrary shift set, its attainment for equal pieces, and its
line form; the bound for decoders that change at single entries (\cref{prop:entry}), apart from a count, and its form
with advice (\cref{cor:entry-advice}); and the twist-class theorem. Not formalized, and proved in the text: the passage from flags to adapted bases; the corollaries
of Theorem~FL that need an algebraically closed field, the structure of $p$-groups or an arithmetic estimate
(\cref{cor:two-terms,cor:commutative,cor:affine,cor:expansion}); the equality case of \cref{prop:swap-floor}; the
copy-rule bound; the count in \cref{prop:entry}; \cref{cor:weyl}; \cref{cor:digit} apart from the abstract lemma it
rests on; the gains for a fixed number of classes over every field
(\cref{prop:two-class,lem:mixed-blowup,cor:every-M}); the bounds for sets
(\cref{prop:sets,prop:order,prop:mutual}) and for whole rows and stages
(\cref{prop:rows,prop:stages,prop:pieces} and \cref{prop:expanding}(a)); and the translation between class-uniform rules and operators $\sum_aA_a\otimes B_a$. Numerical observations appear only in remarks
labelled as numerical and in \cref{sec:small}, and no proof uses them. \Cref{app:leanindex} indexes the declarations and
\cref{app:verification} explains how to check them.

\subsection{Related work}

\paragraph{The series.} \Cref{sec:series} is the account of the earlier papers. The closest of them to the new theorems
are the eighth, which states the coding problem for linear machines as relative rigidity of the transposition
matrix~\cite{EighthPaper2026}; the second and fourth, which exclude locally affine head
motion~\cite{SecondPaper2026,FourthPaper2026}; and the seventh, whose elimination lemmas and Cayley corollary are the
storage-code relatives of \cref{thm:FL,thm:TP}~\cite{SeventhPaper2026}.

\paragraph{Matrix equations with a transpose.} Equations in which $X$ and $X^{\tr}$ both occur have been studied
as $\star$-Sylvester and $\top$-Stein equations. The solution space of $AX+X^{\star}B=0$ is described through the
Kronecker canonical form of a pencil~\cite{DeTeranDopicoGuilleryMontealegreReyes2013}; that of $XA+AX^{\tr}=0$ is
found in~\cite{DeTeranDopico2011} and applied to congruence orbits; unique solvability of $AXB+CX^{\star}D=E$ is
characterized in~\cite{DeTeranIannazzoPoloniRobol2018}, and of $X=AX^{\tr}B+C$ in~\cite{Chiang2013}. Generalized Schur
decompositions are used for solvability and numerical solution~\cite{ChiangChuLin2012}. \Cref{thm:FL} uses
triangular forms for a different purpose, a dimension bound in the singular case with one side unconstrained. We found no statement of \cref{cor:two-terms}, of the
bound for an arbitrary number of terms with one triangularizable side, or of any bound on the rank of the commutation
matrix~\cite{MagnusNeudecker1979} minus a structured sum of Kronecker products. Our search was limited to abstracts and
indexes, and we claim no priority for statements that specialists may regard as folklore.

\paragraph{Commutators and network codes.} \Cref{lem:pigeonhole} is the fact that a commutator has rank at most twice
the rank of either factor; we know no source that states it in the form used here and we give its proof. Transposition
as a multi-commodity network-coding instance goes back to~\cite{AHJKL2006}. Network codes whose local kernels are sums
of cyclic shifts have been studied for multicast~\cite{TangSunLiYangLong2019}; we found nothing on shift-covariant
decoders for transposition.

\subsection{Organization}

\Cref{sec:prelim} fixes the problem and the two models. \Cref{sec:series} is the synthesis and \cref{sec:wall} the
diagnosis. \Cref{sec:order,sec:symmetry} prove the two theorems. \Cref{sec:construction} is the construction side,
\cref{sec:small} the exact computations and \cref{sec:open} the open problems.
